
donc le résultat pour @@M0@@.

Il reste à redescendre jusqu'à @@M1@@. On procède par récurrence descendante, grâce au diagramme commutatif:
@@M2@@
La ligne du bas est exacte car @@M3@@ pour tout @@M4@@ d'après i) déjà prouvé. Par ailleurs, on a:
@@M5@@
La flèche de droite est donc un isomorphisme; celle de gauche en est un d'après l'hypothèse de récurrence; celle du milieu en est donc un également.

\textbf{Proposition 3.6.} \uline{Avec les hypothèses et les notations de} 3.1 \uline{et} 3.3, \uline{on a la relation}:
@@M6@@

Soit @@M7@@; on peut écrire, d'après la formule de projection:
@@M8@@
Il suffit donc de montrer que @@M9@@.

Or 3.5 montre que la cohomologie de @@M10@@ est nulle en degrés @@M11@@, et égale à @@M12@@ en degré @@M13@@. La proposition en résulte aussitôt.

\textbf{Théorème 3.7.} \uline{Soient} @@M14@@ \uline{une immersion fermée régulière}, @@M15@@ \uline{le schéma obtenu en faisant éclater} @@M16@@, \uline{et} @@M17@@, \uline{de sorte qu'on a le diagramme commutatif}
